\section{The 402Pilot Decision Layer}
\label{sec:layer}

402Pilot is a decision layer between an autonomous agent and a programmable payment protocol such as x402. The layer accepts a task and a wallet state, picks a provider, calls the protocol stack to settle the payment, and feeds the realized outcome back into a policy. It does not modify the protocol, never holds funds, and never crosses the settlement boundary. Figure~\ref{fig:architecture} sketches the per-round flow.

\begin{figure}[t]
\centering
\resizebox{\textwidth}{!}{%
\begin{tikzpicture}[
  font=\footnotesize,
  node distance=1.15cm and 1.0cm,
  block/.style={
    draw=black!55,
    rounded corners=2pt,
    align=center,
    minimum height=0.78cm,
    text width=2.15cm,
    inner sep=4pt,
    fill=white
  },
  pilotblock/.style={block, draw=blue!55!black, fill=blue!5},
  extblock/.style={block, draw=black!45, fill=black!7},
  payblock/.style={block, draw=black!45, fill=black!10, dashed},
  arrow/.style={-{Stealth[length=2.0mm]}, line width=0.45pt},
  feedback/.style={-{Stealth[length=2.0mm]}, line width=0.45pt, dashed},
  lab/.style={font=\scriptsize, fill=white, inner sep=1.5pt}
]

\node[extblock] (agent) {Agent\\planner};
\node[pilotblock, right=of agent] (ctx) {Context\\encoder\\$\mathbf{x}_t$};
\node[pilotblock, right=of ctx] (wallet) {Wallet /\\budget\\$\lambda_t,\mathcal{A}_t$};
\node[pilotblock, right=of wallet, text width=2.35cm] (select) {PA-DCT\\\texttt{policy.select}\\$a_t$};
\node[payblock, right=of select, text width=2.3cm] (x402) {Payment\\protocol stack\\(e.g., x402)};
\node[payblock, right=of x402] (provider) {Chosen\\provider\\$a_t$};
\node[extblock, below=1.1cm of x402, text width=2.35cm] (eval) {Outcome\\evaluator\\$q_t,c_t,f_t$};
\node[pilotblock, below=1.1cm of select, text width=2.45cm] (update) {Reward split +\\\texttt{policy.update}\\$u_t=q_t-\nu f_t$};

\begin{scope}[on background layer]
  \node[
    draw=blue!45!black,
    fill=blue!3,
    rounded corners=4pt,
    inner xsep=7pt,
    inner ysep=9pt,
    fit=(ctx) (wallet) (select) (update)
  ] (pilotbox) {};
  \node[
    draw=black!35,
    fill=black!4,
    rounded corners=4pt,
    dashed,
    inner xsep=7pt,
    inner ysep=9pt,
    fit=(x402) (provider)
  ] (paybox) {};
\end{scope}

\node[anchor=north west, font=\scriptsize\bfseries, text=blue!55!black]
  at ([xshift=2pt,yshift=-2pt]pilotbox.north west) {402Pilot decision layer};
\node[anchor=north west, font=\scriptsize\bfseries, text=black!55]
  at ([xshift=2pt,yshift=-2pt]paybox.north west) {delegated settlement / service execution};

\draw[arrow] (agent) -- node[lab, above] {(1) task} (ctx);
\draw[arrow] (ctx) -- node[lab, above] {(2) context} (wallet);
\draw[arrow] (wallet) -- node[lab, above] {(3) budget pressure + feasible arms} (select);
\draw[arrow] (select) -- node[lab, above] {(4) chosen arm} (x402);
\draw[arrow] (x402) -- node[lab, above] {(5) paid call} (provider);
\draw[arrow] (provider.south) |- node[lab, pos=0.28, right] {response} (eval.east);
\draw[arrow] (eval) -- node[lab, below] {(6) observed outcome} (update);
\draw[feedback] (update.north) -- node[lab, right] {(7) update Q/C posteriors} (select.south);
\draw[feedback] (update.west) -| node[lab, pos=0.25, below] {\texttt{record\_spend}$(c_t)$} (wallet.south);

\node[font=\scriptsize, align=center, text=black!65, below=0.18cm of update, text width=4.4cm]
  {Posteriors are updated with utility and realized cost; payment-aware reward is logged for accounting.};

\end{tikzpicture}%
}
\caption{The 402Pilot decision-layer architecture. Per round: (1) task arrives from the agent; (2) the context encoder produces $\mathbf{x}_t$; (3) the wallet exposes $\lambda_t$ and the affordable-arm set; (4) \texttt{policy.select} returns $a_t$; (5) the underlying payment protocol settles and calls the chosen provider (grayed: outside the layer); (6) the evaluator returns $(q_t, c_t, f_t)$; (7) \texttt{policy.update} incorporates the observation. Only steps 2--4 and 7 are inside 402Pilot; settlement is delegated unchanged.}
\label{fig:architecture}
\end{figure}

\paragraph{Two stable interfaces.}
A policy is any object exposing
\begin{align*}
&\mathtt{select}(\mathbf{x}_t,\, \mathcal{A}_t,\, \mathrm{wallet\_state}) \to a_t \in \mathcal{A}_t, \\
&\mathtt{update}(\mathbf{x}_t,\, a_t,\, \ut,\, c_t) \to \emptyset,
\end{align*}
where $\mathcal{A}_t \subseteq \{1, \ldots, K\}$ is the set of providers that fit in the remaining wallet at the current effective prices. Affordability is enforced by the wallet, not the policy: a policy never has to reason about hard budget constraints, only about picking well from a feasible set. The two methods are the only points of contact with the rest of the stack; everything else (priors, posteriors, internal state) is private to the policy.

\paragraph{Update signal: utility, not payment-aware reward.}
A subtle but load-bearing commitment: the layer feeds \emph{utility} $\ut$ to \texttt{policy.update}, not the payment-aware reward $\rt$. This decoupling lets the posterior track \emph{intrinsic provider quality} independently of decision-time budget pressure. If we instead updated on $\rt$, beliefs about a provider would shift purely because the wallet drained, defeating the purpose of having a posterior. Concretely: the same provider should not look worse to the policy late in a budget-tight run than it did early in the run when the wallet was full. Decision-time budget pressure enters only inside \texttt{select}, through $\lamn$, never inside \texttt{update}.

\paragraph{Pluggability across policies and protocols.}
On the policy side, every comparator we evaluate (\S\ref{sec:eval-comparators})---Random, three Always-X variants, BudgetRule, and PA-DCT itself---implements the same two-method interface. New algorithms drop in without touching the rest of the stack. On the protocol side, the layer is agnostic to which programmable payment substrate sits below: today we instantiate above x402~\citep{coinbase2025x402v2}, but the same interfaces apply to AP2~\citep{google2025ap2}, A2A--x402~\citep{a2ax402}, A402~\citep{a402}, or any future quote-and-pay protocol. The protocol stack provides addressing, settlement, and audit trails; 402Pilot provides the decision policy.

\paragraph{What the layer does \emph{not} do.}
Five concerns are explicitly out of scope and are addressed in \S\ref{sec:discuss}:
on-chain settlement (the protocol's job);
capability discovery (handled by x402 V2's Discovery extension or equivalent registries);
production observability (operator dashboards, alerting);
multi-agent coordination (we operate at single-agent scope: one wallet, one policy, one workload stream);
and online safeguards such as drift detection or trust scoring (future work; we discuss directions in \S\ref{sec:discuss}).

\paragraph{What the layer assumes.}
402Pilot makes three assumptions about the surrounding stack. First, the protocol can quote a price for a paid call before the call is made, so the wallet can compute affordability and the policy can include cost in its decision rule. Second, outcomes are observable and scoreable shortly after the call, so the policy can update before the next round (we make this concrete in our benchmark via deterministic per-task evaluators or cached LLM-as-judge scores). Third, providers are non-strategic in the protocol-mechanism sense: they post prices and deliver services rather than playing auction games over bids. All three assumptions hold for x402 V2 and its current peers. Relaxing the third---strategic providers---is the regime studied by recent mechanism-design work~\citep{aamas2026reverseauction}; 402Pilot is complementary to that line, not a substitute.

\paragraph{Implementation footprint.}
The reference implementation is a Python package (\texttt{pilot402}) of approximately [/cc TODO line count] lines, with the policy interface, the wallet, the context encoder, the reward calculator, and the per-seed runtime loop as separate modules; any one of them can be swapped without touching the others. The full benchmark harness, including pre-generated dataset and scenario sweeps, runs end-to-end in approximately [/cc TODO wall-clock estimate; current main sweep takes \~X hours on a laptop] and is released alongside this paper.
